\subsection{Системы линейных алгебраических уравнений. Основные определения. Теорема Крамера}

\begin{definition}
Системой линейных алгебраических уравнений (СЛАУ) называется система вида:
\[
\begin{cases}
a_{11}x_1 + a_{12}x_2 + \dots + a_{1n}x_n = b_1 \\
a_{21}x_1 + a_{22}x_2 + \dots + a_{2n}x_n = b_2 \\
\dots \\
a_{m1}x_1 + a_{m2}x_2 + \dots + a_{mn}x_n = b_m
\end{cases}
\]
где $a_{ij}$ — коэффициенты системы, $b_i$ — свободные члены, $x_j$ — неизвестные.
\end{definition}

Для СЛАУ вводятся следующие матричные обозначения:
\begin{itemize}
    \item \textbf{Основная матрица системы} $A$ размера $m \times n$, составленная из коэффициентов $a_{ij}$.
    \item \textbf{Столбец неизвестных} $X = (x_1, \dots, x_n)^T$.
    \item \textbf{Столбец свободных членов} $B = (b_1, \dots, b_m)^T$.
    \item \textbf{Расширенная матрица системы} $\tilde{A}$ (или $A$ с чертой), полученная приписыванием столбца $B$ к матрице $A$.
\end{itemize}

В матричной форме система записывается как:
\[ AX = B \]

\begin{definition}
Система называется \textbf{однородной}, если столбец свободных членов $B = \mathbf{0}$. В противном случае система называется \textbf{неоднородной}.
\end{definition}

\begin{definition}
\textbf{Решением системы} называется упорядоченный набор чисел $(\gamma_1, \dots, \gamma_n)$, который при подстановке в систему вместо неизвестных превращает каждое уравнение в верное тождество.
\end{definition}

Классификация систем по количеству решений:
\begin{enumerate}
    \item \textbf{Совместная} — имеет хотя бы одно решение.
    \begin{itemize}
        \item \textbf{Определенная} — имеет ровно одно решение.
        \item \textbf{Неопределенная} — имеет более одного решения.
    \end{itemize}
    \item \textbf{Несовместная} — не имеет ни одного решения.
\end{enumerate}

\begin{theorem}[Крамера]
Пусть дана квадратная СЛАУ ($m = n$). Если определитель основной матрицы $\Delta = \det A \neq 0$, то система имеет единственное решение, которое находится по формулам:
\[ x_1 = \frac{\Delta_1}{\Delta}, \dots, x_n = \frac{\Delta_n}{\Delta} \]
где $\Delta_k$ — определитель, полученный из $\Delta$ заменой $k$-го столбца столбцом свободных членов $B$.
\end{theorem}

\begin{tcolorbox}[title=Доказательство, breakable]
\textbf{1. Существование и единственность.}
Запишем систему в матричном виде: $AX = B$. Так как $\det A \neq 0$, матрица $A$ является невырожденной, и для неё существует единственная обратная матрица $A^{-1}$. Умножим обе части матричного равенства слева на $A^{-1}$:
\[ A^{-1}(AX) = A^{-1}B \implies (A^{-1}A)X = A^{-1}B \implies EX = A^{-1}B \implies X = A^{-1}B \]
Поскольку обратная матрица и результат умножения определены однозначно, решение существует и оно единственно.

\textbf{2. Вывод формул Крамера.}
Воспользуемся формулой для обратной матрицы через присоединенную (матрицу алгебраических дополнений):
\[ A^{-1} = \frac{1}{\Delta} \begin{pmatrix} A_{11} & A_{21} & \dots & A_{n1} \\ A_{12} & A_{22} & \dots & A_{n2} \\ \dots & \dots & \dots & \dots \\ A_{1n} & A_{2n} & \dots & A_{nn} \end{pmatrix} \]
Тогда столбец решений $X$ вычисляется как:
\[ \begin{pmatrix} x_1 \\ \dots \\ x_k \\ \dots \\ x_n \end{pmatrix} = \frac{1}{\Delta} \begin{pmatrix} A_{11} & A_{21} & \dots & A_{n1} \\ \dots & \dots & \dots & \dots \\ A_{1k} & A_{2k} & \dots & A_{nk} \\ \dots & \dots & \dots & \dots \\ A_{1n} & A_{2n} & \dots & A_{nn} \end{pmatrix} \begin{pmatrix} b_1 \\ b_2 \\ \dots \\ b_n \end{pmatrix} \]
Для произвольного $k$-го неизвестного получаем:
\[ x_k = \frac{1}{\Delta} \sum_{i=1}^n A_{ik} b_i \]
Сумма в правой части $\sum_{i=1}^n b_i A_{ik}$ представляет собой разложение определителя $\Delta_k$ по $k$-му столбцу (в котором стоят элементы $b_i$, а соответствующие им алгебраические дополнения $A_{ik}$ такие же, как в определителе $\Delta$, так как $k$-й столбец при их вычислении вычеркивается).
Следовательно:
\[ x_k = \frac{\Delta_k}{\Delta} \]
Теорема доказана.
\end{tcolorbox}

\begin{tcolorbox}[title=Источники, colback=gray!10]
\begin{itemize}
  \item Лекция №\,2, тема «Системы линейных алгебраических уравнений. Основные понятия», временной отрезок 52:00--61:00
  \item Лекция №\,2, тема «Теорема Крамера», временной отрезок 62:00--68:00
  \item PDF «Линейная алгебра\_compressed (1).pdf», разделы 1 и 2 по СЛАУ, стр. 7--10
  \item PDF «Вопросы», вопрос №\,4
\end{itemize}
\end{tcolorbox}
